
\subsection{The monoidal product from the wedge product}

\begin{definition}[Monoidal product]\label{def:monoidal-prod}
In the category $\mathcal{C}(S,\kappa)$, the
\emph{monoidal product} of two objects $A$ and $B$ is their
joint discrimination profile: the object $A \otimes B$ whose
$\tau/\sigma$ profile under each discriminator $d \in \kappa$
is determined by the simultaneous application of $d$ to both
$A$ and $B$.

For morphisms, the monoidal product of $f: A \to B$
(witnessed by $d_f$) and $g: C \to D$ (witnessed by $d_g$)
is $f \otimes g: A \otimes C \to B \otimes D$, witnessed by
the wedge product $d_f \wedge d_g$ --- the joint, parallel
application of both discriminations.
\end{definition}

\begin{proposition}[Symmetric monoidal structure]
\label{prop:sym-monoidal}
$(\mathcal{C}(S,\kappa), \otimes)$ is a symmetric monoidal
category. Specifically:
\begin{enumerate}[label=(\roman*),nosep]
  \item \textbf{Associativity}: $(A \otimes B) \otimes C
    \cong A \otimes (B \otimes C)$, inherited from
    associativity of the wedge product.
  \item \textbf{Unit}: the monoidal unit $I$ is the
    degenerate object (the grade-0 element), representing
    the trivial discrimination.
  \item \textbf{Symmetry}: $A \otimes B \cong B \otimes A$,
    because $d_f \wedge d_g = d_g \wedge d_f$ over $\GF(2)$.
\end{enumerate}
\end{proposition}

\begin{proof}
The wedge product on $\bigwedge(\GF(2)^n)$ is associative,
unital (with unit $1 \in \bigwedge^0 V$), and commutative
over $\GF(2)$. The monoidal product is defined via the
wedge product. Therefore $(\mathcal{C}(S,\kappa), \otimes)$
inherits symmetric monoidal structure from the exterior
algebra.
\end{proof}

\begin{remark}
The monoidal product (parallel composition) and categorical
composition (sequential composition) are both mediated by
the wedge product, but they differ in the $\tau/\sigma$
tracking. Sequential composition
(Definition \ref{def:composition}) requires the intermediate
object to appear on opposite sides ($\sigma$ of the first
discriminator, $\tau$ of the second). Parallel composition
imposes no such matching condition: the two discriminations
are applied independently. The algebra is the same; the
$\tau/\sigma$ choreography differs.
\end{remark}

\subsection{Delooping as rotation}\label{sec:delooping}

\begin{theorem}[Delooping equivalence]\label{thm:deloop}\footnote{Axiom class \textsf{A} $\cdot$ depth 5 $\cdot$ \S\ref{app:deloop-closure}.  \emph{Warrant}: standard delooping; a one-object 2-category with $\wedge$ as monoidal product. Finite-$\kappa$: yes. Standard (one-object 2-category = monoidal category).}
Let $\mathcal{C}(S,\kappa)$ be the $n$-category
viewed from the $\kappa$-perspective. Rotating to the
$\sigma$-perspective (Definition \ref{def:triangle},
Proposition \ref{prop:perspectives}) produces a monoidal
$(n-1)$-category:
\begin{itemize}[nosep]
  \item What were 1-morphisms (discriminators in $\kappa$)
    become \emph{objects} (elements of the population of
    the $\sigma$-regime).
  \item What was the wedge product of 1-morphisms (their
    composition/tensor) becomes the \emph{monoidal product
    of objects}.
  \item What were 2-morphisms (relationships between
    discriminators) become 1-morphisms in the new category.
\end{itemize}
This is delooping: a 2-category with one object is
equivalently a monoidal 1-category. The rotation of
the triangle performs delooping as a geometric operation.
\end{theorem}

\begin{proof}
In the $\kappa$-perspective, the population is $S$ and the
discriminators are elements of $\kappa$. In the
$\sigma$-perspective (equation \ref{eq:tri-sigma}), the
population is $\kappa$ (the discriminators are now the
things being discriminated) and the regime is $\sigma$.

An object in the $\sigma$-perspective is a
$\sigma$-equivalence class of elements of $\kappa$, i.e.\
a class of discriminators. A 1-morphism in the
$\sigma$-perspective is a discriminator $d_\sigma \in \sigma$
that separates discriminators of $\kappa$ into $\tau$ and
$\kappa$-masks (using the rotated labeling).

The wedge product of two discriminators
$d_1, d_2 \in \kappa$ --- which was a 2-cell
$d_1 \wedge d_2$ in the $\kappa$-perspective --- is now
the monoidal product $d_1 \otimes d_2$ of two objects in
the $\sigma$-perspective. The algebraic operation is
identical; only the categorical interpretation changes.

This is precisely the Baez--Dolan delooping: an
$(n+1)$-category with one object is equivalently a
monoidal $n$-category. The single object in the
$(n+1)$-category corresponds to the entire population
$S$ (which collapses to a point when viewed from the
meta-level). The morphisms of the $(n+1)$-category become
the objects of the monoidal $n$-category. The rotation
of the triangle performs this identification without
constructing any new structure.
\end{proof}

\begin{corollary}[Delooping collapse over $\GF(2)$]
\label{cor:iterated-deloop}
Over $\GF(2)$, the braided/symmetric distinction in the
Baez--Dolan periodic table \emph{collapses}: since
$d \wedge d = 0$ and $-1 = 1$, the braiding
$\beta: A \otimes B \to B \otimes A$ satisfies
$\beta^2 = \mathrm{id}$ immediately. The braiding is an
involution in one step, not two. Therefore the second and
third deloopings produce the same structure:
\begin{itemize}[nosep]
  \item $\kappa$-perspective: $n$-category.
  \item $\sigma$-perspective: monoidal $(n{-}1)$-category.
  \item $\tau$-perspective: symmetric monoidal
    $(n{-}2)$-category (not merely braided ---
    the braided and symmetric rows are identified over
    $\GF(2)$).
\end{itemize}
The exterior algebra, being commutative at all levels,
lives permanently on the ``symmetric monoidal'' row of
the periodic table. The rows below it (braided,
monoidal) are invisible to $\bigwedge(\GF(2)^n)$.
\end{corollary}

\begin{proof}
The braiding $\beta_{A,B}: A \otimes B \to B \otimes A$
is the swap of tensor factors. Over a field of
characteristic $\neq 2$, the sign $(-1)^{|A||B|}$ in the
graded commutativity relation distinguishes the braiding
from the identity, and the hexagon axioms produce
non-trivial structure. Over $\GF(2)$,
$(-1)^{|A||B|} = 1$ for all grades, so
$\beta = \mathrm{id}$ as a natural transformation. The
braiding is trivially symmetric. No second delooping is
needed to reach the symmetric row; the first delooping
already lands there.
\end{proof}

\subsection{Cayley--Dickson recovery of the full periodic
  table}\label{sec:cayley-dickson}

The collapse of the periodic table over $\GF(2)$ is a
genuine structural phenomenon, not an error. But the
full periodic table --- with distinct braided and
monoidal rows --- can be recovered using the
\emph{Cayley--Dickson construction} applied to the
subobject classifier $\Omega = \GF(2)^2$.

\begin{definition}[Swap conjugation]\label{def:swap-conj}
The \emph{swap conjugation} on $\GF(2)^2$ is the map
$(a, b)^* = (b, a)$. This is the negation operator
$\neg$ on the subobject classifier (Remark after
Theorem \ref{thm:topos}): it swaps the $\tau$ and
$\sigma$ components. Over $\GF(2)$, this is a
non-trivial involution even though $-1 = 1$: the
\emph{sign} is dead, but the \emph{swap} is alive.
\end{definition}

\begin{definition}[Cayley--Dickson tower over $\GF(2)^2$]
\label{def:cayley-dickson}
The Cayley--Dickson construction takes an algebra $A$
with conjugation $*$ and produces a doubled algebra
$\mathrm{CD}(A) = A \times A$ with multiplication
$(a, b)(c, d) = (ac + d^* b,\; da + bc^*)$ and
conjugation $(a, b)^* = (a^*, -b) = (a^*, b)$ (since
$-1 = 1$ over $\GF(2)$). Iterating:
\begin{itemize}[nosep]
  \item \textbf{Step 0}: $\GF(2)$ --- scalars. The
    monoidal unit. Commutative, associative.
  \item \textbf{Step 1}: $\GF(2)^2$ with swap conjugation
    --- the four-cell structure. Commutative, associative.
    Corresponds to \emph{symmetric monoidal}.
  \item \textbf{Step 2}: $(\GF(2)^2)^2 = \GF(2)^4$ ---
    \emph{non-commutative}, still associative. The swap
    conjugation in the Cayley--Dickson product breaks
    commutativity: left-multiplication and
    right-multiplication are distinguishable. Corresponds
    to \emph{braided monoidal}: the braiding carries
    non-trivial structure because the algebra distinguishes
    left from right.
  \item \textbf{Step 3}:
    $((\GF(2)^2)^2)^2 = \GF(2)^8$ ---
    non-commutative, \emph{non-associative}. Corresponds
    to \emph{monoidal} (no braiding): the coherence
    conditions are non-trivial because associativity fails.
\end{itemize}
\end{definition}

\begin{theorem}[The periodic table is two-dimensional]\footnote{Axiom class \textsf{A}. \emph{Warrant}: in-universe recognizer construction via $\sigma$-rotation. \emph{Hexagon axioms}: hold by strictness (same argument as G12). Zero divisors at step 2 are residues ($\tau \cdot \sigma = 0$), not obstructions. See Appendix \ref{app:cd-recognizer}. Finite-$\kappa$: yes.}
\label{prop:periodic-2d}
The full Baez--Dolan periodic table over $\GF(2)$ is
the product of two independent structures:
\begin{itemize}[nosep]
  \item The \emph{grading} of $\bigwedge(\GF(2)^n)$ gives
    the categorical dimension $n$ (the vertical axis).
  \item The \emph{Cayley--Dickson level} over $\GF(2)^2$
    gives the monoidal depth $k$ (the horizontal axis):
    $k = 0$ is the exterior algebra itself (symmetric),
    $k = 1$ is the Cayley--Dickson doubling
    (braided/non-commutative), $k = 2$ is the double
    doubling (monoidal/non-associative).
\end{itemize}
The two coordinates are algebraically independent: the
grading measures how many discriminators have been
composed, while the Cayley--Dickson level measures how
many layers of swap-conjugation structure have been
introduced.
\end{theorem}

\begin{proof}
The exterior algebra $\bigwedge(\GF(2)^n)$ is commutative
at all grades (Proposition \ref{prop:commutative}), so
it cannot distinguish braided from symmetric monoidal
structure. The Cayley--Dickson construction recovers the
distinction: at CD step 1, the algebra is commutative
(symmetric); at CD step 2, non-commutative (braided); at
CD step 3, non-associative (merely monoidal). Each CD step
peels off one layer of algebraic regularity, corresponding
to one row of the periodic table.

\emph{Caveat}: the Baez--Dolan correspondence was
established for the CD tower over $\mathbb{R}$
(reals $\to$ complexes $\to$ quaternions $\to$
octonions). The tower over $\GF(2)^2$ is a different
object, and the correspondence does not transfer
automatically. This section is verified via in-universe recognizer
construction (Appendix \ref{app:cd-recognizer}), pending
verification of the hexagon axioms.

The grading and the CD level are independent because the
grading operates on the \emph{number} of composed
discriminators while the CD level operates on the
\emph{conjugation structure} of the base field. One can
compose arbitrarily many discriminators (increasing the
grade) at any fixed CD level, and one can double the
conjugation structure (increasing the CD level) at any
fixed grade.
\end{proof}  % end of periodic table proof

\begin{remark}[The swap conjugation is already present]
The Cayley--Dickson construction requires no new
structure. The swap conjugation $(a,b)^* = (b,a)$ is the
negation operator $\neg$ on $\Omega$, already present in
the topos structure (Section \ref{sec:topos}). Each
Cayley--Dickson doubling is the iterated application of
this negation as a conjugation, doubling the algebra at
each step. Each doubling is interpretable as a
$\kappa$-evolution that adds a \emph{conjugation dimension}:
a dimension that tracks not just $\tau/\sigma$ membership
but the \emph{orientation} of the discrimination.
\end{remark}

\begin{remark}[The periodic table as the toroid
  $\times$ the CD tower]
If the Cayley--Dickson correspondence holds over
$\GF(2)^2$ (which depends on the open hexagon axiom
verification of Remark 8.11 below), then
the toroid provides the categorical
dimension and the symmetric monoidal column, and the
Cayley--Dickson tower provides the remaining columns.
The full periodic table would be the product:
\[
  \text{Periodic table} \stackrel{?}{=} \bigwedge(\GF(2)^n)
  \times \mathrm{CD}^k(\GF(2)^2).
\]
The toroid's geometry gives the $n$-coordinate; the
Cayley--Dickson level gives the $k$-coordinate.
This product structure is conditional on the CD tower
preserving the expected algebraic properties at each
step; the zero-divisor issue at step 2 (Appendix
\S\ref{app:cd-table}) may cause earlier degeneration
than the classical tower over $\mathbb{R}$.
(Delooping and closure verification: \S\ref{app:deloop-closure}.)
\end{remark}

\begin{remark}[Open verification: hexagon axioms at CD step 2]
The correspondence between Cayley--Dickson algebraic
properties and categorical structures ---
non-commutativity at CD step 2 corresponding to braided
monoidal structure, non-associativity at CD step 3
corresponding to merely monoidal structure --- follows
the standard pattern identified by Baez and Dolan.
However, explicitly constructing the braiding at CD
step 2 (as a natural isomorphism
$\beta_{A,B}: A \otimes B \to B \otimes A$ in the
$\GF(2)^4$ algebra), verifying the hexagon axioms, and
confirming that the symmetry condition
$\beta_{B,A} \circ \beta_{A,B} = \mathrm{id}$ fails
(i.e.\ that the braiding is genuinely non-symmetric)
remains an open verification. Additionally, the
Cayley--Dickson tower over $\GF(2)^2$ may degenerate
faster than the classical tower over $\mathbb{R}$
because $\GF(2)^2$ has zero divisors
($(1,0) \cdot (0,1) = (0,0)$), potentially losing
alternativity at step 2 rather than step 3. This
deserves explicit computation. (A partial multiplication table confirming non-commutativity and zero divisors is in \S\ref{app:cd-table}.)
\end{remark}

\subsection{Internal hom from the triangle identity}

\begin{definition}[Internal hom]\label{def:internal-hom}
For objects $A$ and $B$ in $\mathcal{C}(S,\kappa)$, the
\emph{internal hom} $[A, B]$ is the collection of all
discriminators $d \in \kappa$ such that $\tau_d(A) = 1$
and $\sigma_d(B) = 1$. This collection is itself an
object in the $\sigma$-perspective: it is a sub-population
of $\kappa$ (viewed as a population under the
$\sigma$-regime).
\end{definition}

\begin{theorem}[Monoidal closure]\label{thm:closure}\footnote{Axiom class \textsf{A} $\cdot$ depth 4 $\cdot$ \S\ref{app:deloop-closure}.  \emph{Warrant}: tensor-hom adjunction from the triangle identity. $[B,C]$ is the set of discriminators witnessing $B \to C$. Finite-$\kappa$: yes. Standard (tensor-hom adjunction).}
The symmetric monoidal category
$(\mathcal{C}(S,\kappa), \otimes)$ is closed: there
exists a natural bijection
\begin{equation}\label{eq:closure}
  \operatorname{Hom}(A \otimes B, C) \cong
  \operatorname{Hom}(A, [B, C])
\end{equation}
for all objects $A, B, C$.
\end{theorem}

\begin{proof}
The triangle identity (Definition \ref{def:triangle})
states that any two of $\{\kappa, \sigma, \tau\}$
determine the third. We apply this at the level of
morphisms.

A morphism $A \otimes B \to C$ is a discriminator $d$
with $\tau_d(A \otimes B) = 1$ and $\sigma_d(C) = 1$.
Since $A \otimes B$ is the joint discrimination profile
of $A$ and $B$, this means $\tau_d(A) = 1$,
$\tau_d(B) = 1$, and $\sigma_d(C) = 1$.

A morphism $A \to [B,C]$ is a discriminator $d'$
with $\tau_{d'}(A) = 1$ and $\sigma_{d'}([B,C]) = 1$.
Now $[B,C]$ is the set of discriminators $e$ with
$\tau_e(B) = 1$ and $\sigma_e(C) = 1$. So
$\sigma_{d'}([B,C]) = 1$ means $d'$ classifies $A$ on
the $\tau$ side and identifies a discriminator $e$
that sends $B$ to $\tau$ and $C$ to $\sigma$.

The bijection is: given $d$ witnessing $A \otimes B \to C$,
set $d' = d$ and $e = d$ (since $d$ already has
$\tau_d(A) = 1$, $\tau_d(B) = 1$, $\sigma_d(C) = 1$).
Conversely, given $d'$ and $e$, the discriminator $d$
that jointly witnesses both is their wedge product in the
$\sigma$-regime.

This bijection is natural in $A$ because the $\tau/\sigma$
assignments are functorial (preserved by $\kappa$-deformations).
The triangle identity guarantees that the three-way
relationship $\kappa = \sigma + \tau$ --- here
instantiated as the relationship among source, target,
and hom-space --- is consistent and lossless.
\end{proof}

\begin{remark}[Closure is the triangle identity]
The adjunction (\ref{eq:closure}) is the triangle identity
applied at the morphism level. The three components ---
$A \otimes B$, $C$, and $[B,C]$ --- stand in the same
triadic relationship as $\kappa$, $\sigma$, and $\tau$:
any two determine the third. Monoidal closure is not an
additional property that the category may or may not
possess. It is a direct expression of the $\GF(2)$
toroid's geometry. The category is closed because the
triangle is closed.
\end{remark}

\begin{definition}[Evaluation morphism]\label{def:eval}
The \emph{evaluation morphism}
$\mathrm{ev}: [A, B] \otimes A \to B$ is the morphism
witnessed by the application of a discriminator in $[A,B]$
to the element $A$ in the population. This is the interface
between the $\sigma$-regime (which manages the hom-space)
and the $\kappa$-regime (which manages the population).
The evaluation morphism is the bridge between perspectives:
it takes an object from the $\sigma$-perspective (a
discriminator) and applies it in the $\kappa$-perspective
(to a population element), producing a result in the
original category.
\end{definition}

\begin{remark}[Monoidal structure and $n$-categorical
  structure]
The $n$-categorical hierarchy and the symmetric monoidal
column of the periodic table are equivalent under
rotation of the $\GF(2)$ toroid:
\begin{itemize}[nosep]
  \item Fixing a vertex and reading off the grading gives
    $n$-categorical structure.
  \item Rotating and collapsing a level gives symmetric
    monoidal structure (Corollary \ref{cor:iterated-deloop}).
  \item The closure (internal hom) is the third vertex of
    the triangle determined by the other two.
\end{itemize}
The full monoidal hierarchy --- including the braided and
merely-monoidal rows --- requires the Cayley--Dickson
tower (Section \ref{sec:cayley-dickson}), which is built
from the swap conjugation already present in the
subobject classifier. The complete periodic table is the
product of the exterior algebra (categorical dimension)
and the Cayley--Dickson level (monoidal depth).
\end{remark}

